% ECONOMICS_PROBLEM: {"id": "C73-11", "title": "Characterizing universal convergence of fictitious play", "jel_code": "C73", "source_id": "OP-0136"}
% FIRST_POSTED: 2026-10-11
% ATTEMPT_STATUS: solved
% ENTRY_KIND: research_attempt
% SUBMISSION: {"kind": "research_attempt", "category": "economics", "problem_id": "C73-11", "model": {"provider": "Anthropic", "version": "Claude Fable 5.1 (claude-fable-5-1)", "reasoning": "Claude Code agent session, default effort", "generated_on": "2026-10-11"}, "other_models": "None for the Lean development; the hand proof that is formalized was generated by GPT Codex (pull request #138, manuscript linked below)", "tools": "Lean 4.34.1 with Mathlib d13f23b723b8a846827a245b89c10fc7d3f11612 and lake; compilation on a separate Linux server", "target": "The exact C73-11 target as stated by the hand proof: necessary and sufficient conditions on the finite payoff arrays for simultaneous fictitious play to approach the Nash set from every initial profile and under every legal tie choice, with all quantifiers retained", "scope": "Theorem 1 ((i) <=> (ii) <=> (iii)) and Proposition 3 (the class is coanalytic) are formalized in full; the Tarski quantifier-elimination compilation of the admission tests (a presentational rewriting) and the rank-omega example are not formalized", "human_contributions": "None; the submitter specified the task and submitted the result, with no mathematical edits", "verification": "Lean kernel check of every statement with only the standard axioms (propext, Classical.choice, Quot.sound); model self-audit of the definitions against the manuscript; no independent human proof verification", "reproducibility": "Public Lean repository pinned to a commit, with build commands and the actual checking outcome reported; the interactive agent-session prompts are not published"}
% FULL_SOLUTION_SCOPE: {"target": "The entire original C73-11 target as formulated by the hand proof (Theorem 1 of the manuscript), including the coanalytic bound (Proposition 3)", "result": "main_theorem in FPlay/Main.lean: CFP u <-> CondO u, CondO u <-> TreeWF u, TreeWF u <-> no infinite branch; analyticSet_compl_CFP in FPlay/Coanalytic.lean", "coverage": "Every finite game (any finite player set and nonempty finite action sets), all initial profiles including boundary faces, all legal tie choices, the exact Nash set with Euclidean distance; no existence theorem for Nash equilibria is used", "dependencies": "Mathlib only (d13f23b7); no additional axioms; the hand proof's use of Tarski quantifier elimination is a presentational rewriting and is not formalized", "human_verification": "None", "unverified": "The transcription of the manuscript's definitions into Lean was checked by the model only; the hand proof itself is an LLM claim pending community expert review; the rank-omega example is not formalized", "novelty": "Not checked beyond the manuscript's own references; the survey of Li et al. (2023) lists the characterization as open"}
\documentclass[11pt]{article}
\usepackage[margin=1in]{geometry}
\usepackage{amsmath,amssymb,amsthm}
\usepackage[hidelinks]{hyperref}
\newtheorem{theorem}{Theorem}
\newtheorem{proposition}[theorem]{Proposition}
\theoremstyle{definition}
\newtheorem{definition}[theorem]{Definition}
\newcommand{\N}{\mathbb N}
\newcommand{\R}{\mathbb R}
\DeclareMathOperator{\NE}{NE}
\DeclareMathOperator{\dist}{dist}
\title{A complete machine-checked proof of the payoff-array characterization of universal convergence of fictitious play}
\author{Claude Fable 5.1 (Claude Code), submitted by Ji Ho Bae}
\date{October 11, 2026}
\begin{document}
\maketitle
\begin{abstract}
The hand proof for this problem characterizes the finite normal-form games in which simultaneous fictitious play approaches the Nash set from every initial profile and under every legal tie choice, by a countable tree built from finite polynomial-sign tests on the payoff entries: the game is in the class if and only if the tree admits a strictly decreasing countable-ordinal labeling, equivalently is well founded. This note reports a complete Lean~4 formalization of that proof against Mathlib: all three assertions of the theorem are proved equivalent for every finite game, and the coanalytic upper bound on the class is proved as well. Convergence to the Nash set is formulated literally, so that no existence theorem for Nash equilibria is needed anywhere. The development has no \texttt{sorry}, no warnings, assumes nothing beyond Mathlib, and uses only the standard axioms.
\end{abstract}

\section{The statement}

Let $G=(N,(S_i),(u_i))$ be a finite normal-form game with $n=|N|\ge1$ players and nonempty finite action sets; $A=\prod_iS_i$, $X=\prod_i\Delta(S_i)$, payoffs on $X$ are multilinear extensions and $u_i(b,x_{-i})$ is the payoff of the pure deviation $b$. From any $x^1\in X$, simultaneous fictitious play chooses at every $t\ge1$ a profile of pure best replies $a^{t+1}$ to $x^t$ (every tie choice is legal) and sets $x^{t+1}_i=(tx^t_i+e_{a^{t+1}_i})/(t+1)$. The class $\mathcal C_{FP}$ consists of the games for which $\dist(x^t,\NE(G))\to0$ for every $x^1$ and every legal sequence, with the Euclidean distance.

A \emph{record} consists of a threshold index $m\ge1$, a joint-action word $v=(a^2,\dots,a^T)$, marked times $1\le t_1<\dots<t_r=T$ and witnesses $(i_\ell,b_\ell)$. With $Z^t_{i,b}(\zeta)=\zeta_{i,b}+\#\{q\le t:a^q_i=b\}$, $P^t_{i,b}=\sum_{s_{-i}}u_i(b,s_{-i})\prod_{j\ne i}Z^t_{j,s_j}$ and $H^t_{i,b}=tP^t_{i,b}-\sum_dZ^t_{i,d}P^t_{i,d}$, the feasible set $C_q(u)\subseteq X$ is cut out by the legality inequalities (L) $P^t_{i,a^{t+1}_i}\ge P^t_{i,b}$ for $t<T$ and the marked inequalities (B) $mH^{t_\ell}_{i_\ell,b_\ell}\ge t_\ell^{\,n}$. A record is admitted when $C_q(u)\ne\varnothing$. The tree $\mathcal T(u)$ has a root and one vertex per admitted record; the parent of a record deletes its last mark and witness and truncates the word accordingly, so that an edge appends one nonempty action block ending at a new mark (only the first block may be empty).

\begin{theorem}\label{main}
For every finite game the following are equivalent:
\textup{(i)} $G\in\mathcal C_{FP}$;
\textup{(ii)} condition \textup{(O)}: there are a countable ordinal $\alpha$ and $\rho:\mathcal T(u)\to\alpha$, root included, strictly decreasing from every vertex to each of its children;
\textup{(iii)} $\mathcal T(u)$ is well founded, i.e.\ has no infinite branch.
\end{theorem}

\begin{proposition}\label{coan}
Each $Q_q=\{u:C_q(u)\ne\varnothing\}$ is closed in $\R^{n|A|}$, and $\mathcal C_{FP}$ is coanalytic.
\end{proposition}

\section{What is verified}

The Lean development is public at a fixed commit \cite{Lean}. Its objects are the ones above: \texttt{Prof} and \texttt{IsMixed} (mixed profiles), \texttt{mult} and \texttt{dev} (multilinear payoffs and pure deviations), \texttt{NE} and \texttt{resid} (the Nash set and the residual), \texttt{Legal} (the rule (fp)), \texttt{ConvNE} and \texttt{CFP}, \texttt{Record}, \texttt{Feas} ($C_q(u)$), \texttt{Admitted}, the vertex set \texttt{Vertex} and child relation \texttt{TChild} of $\mathcal T(u)$, \texttt{TreeWF} (iii), \texttt{HasBranch} (an infinite branch from the root), and \texttt{CondO} (ii), the latter through \texttt{HasCountableRank}: a labeling by ordinals below a countable ordinal, strictly decreasing along the child relation.

\subsection*{The theorem}
\begin{verbatim}
theorem main_theorem (u : Payoffs N S) :
    (CFP u <-> CondO u) /\ (CondO u <-> TreeWF u) /\ (TreeWF u <-> not HasBranch u)
\end{verbatim}
Unicode is transliterated here (\texttt{<->}, \texttt{/\textbackslash}, \texttt{not}, \texttt{forall}, \texttt{->}). The statement has no hypothesis beyond the finiteness and nonemptiness of the player and action sets. The two directions of (i)$\Leftrightarrow$(iii) are exported separately (\texttt{hasChain\_of\_not\_CFP}, \texttt{not\_CFP\_of\_hasChain}), as is (ii)$\Leftrightarrow$(iii) (\texttt{condO\_iff\_treeWF}).

\subsection*{The coanalytic bound}
\begin{verbatim}
theorem isClosed_Qset (q : Record N S) : IsClosed (Qset q)

theorem analyticSet_compl_CFP : MeasureTheory.AnalyticSet {u : Payoffs N S | not CFP u}
\end{verbatim}
Here \texttt{Payoffs N S} is the Polish space $\R^{n|A|}$ of payoff arrays and \texttt{AnalyticSet} is Mathlib's notion; the complement of the class is analytic, so the class is coanalytic.

\subsection*{Convergence to the Nash set}
\texttt{ConvNE u x} states that for every $\varepsilon>0$, eventually some Nash profile is within Euclidean distance $\varepsilon$ of $x^t$. This is the literal statement $\dist(x^t,\NE(G))\to0$ under the convention $\dist(x,\varnothing)=+\infty$, and with it both directions of Lemma~2 of the hand proof hold without any existence theorem for Nash equilibria; the hand proof invokes Nash's theorem only at that point. When the Nash set is nonempty, \texttt{ConvNE} is the convergence of Mathlib's \texttt{Metric.infDist} to zero (\texttt{convNE\_iff\_infDist}).

\section{Architecture of the formal proof}

\begin{itemize}
\item \emph{Lemma 1} (Section 4). For the trajectory from $\zeta$ along a word, $tx^t=Z^t$, $u_i(b,x^t_{-i})=P^t_{i,b}/t^{n-1}$ and $u_i(b,x^t_{-i})-u_i(x^t)=H^t_{i,b}/t^n$ are proved from the homogeneity of the multilinear extension and the identity $u_i(x)=\sum_dx_{i,d}u_i(d,x_{-i})$; hence (L) is exactly legality of the word and (B) is exactly a gain of at least $1/m$ (\texttt{isBR\_iff\_L}, \texttt{gain\_iff\_B}).
\item \emph{Lemma 2} (Section 4). The residual is continuous, nonnegative on $X$ and zero exactly on the Nash set; distance convergence implies residual convergence by uniform continuity on the compact $X$, and the converse by extracting a convergent subsequence from a subsequence staying $\varepsilon$ away from the Nash set (\texttt{convNE\_iff\_tendsto\_resid}). Failure is one fixed threshold $1/m$ exceeded infinitely often (\texttt{not\_convNE\_iff}).
\item \emph{Records} (Sections 2--3). A record is encoded as $m$ with a list of blocks, each block the actions between two consecutive marks and the witness at the new mark; the word, the last time $T$ and the marks are read off the blocks. The feasible set of a child is contained in the feasible set of its parent, since every constraint of the parent is among those of the child (\texttt{feas\_subset\_of\_child}); feasible sets are closed and compact.
\item \emph{Branches} (Section 6). From a failing legal orbit, Lemma~2 gives $m$ and strictly increasing marked times with witnesses; the records of the marked prefixes are admitted because the actual initial profile satisfies all their constraints (Lemma~1), and they form an infinite chain. Conversely an infinite chain of admitted records has nested nonempty compact feasible sets, hence a common initial profile; its orbit along the union word is legal and has residual at least $1/m$ at every marked time. Chains and branches from the root are interchanged by prepending the truncations of the first record (\texttt{hasBranch\_iff\_hasChain}).
\item \emph{Ordinals} (Section 7). For a relation on a countable type, well-foundedness is equivalent to a strictly decreasing labeling by a countable ordinal: the rank function of Mathlib's well-founded recursion has countable values because every vertex has countably many children, by the regularity of $\aleph_1$ (\texttt{wellFounded\_iff\_hasCountableRank}); well-foundedness is also the absence of infinite descending sequences (\texttt{wellFounded\_iff\_no\_chain}).
\item \emph{Coanalyticity} (Section 8). The constraints are jointly continuous in $(u,\zeta)$, so $Q_q$ is the projection of a closed set along the compact $X$ and is closed. With the records enumerated by $\N$, the set of pairs $(u,p)\in\R^{n|A|}\times\N^{\N}$ naming an infinite branch from the root is closed; its projection is the complement of the class by Theorem~\ref{main}, and the continuous image of a closed subset of a Polish space is analytic.
\end{itemize}

\section{Scope}

The formalization follows the hand proof. The tree is defined by the feasibility $C_q(u)\ne\varnothing$ of the finite polynomial system; the hand proof additionally compiles each admission test into a quantifier-free formula $\Theta_q(u)$ by Tarski's quantifier elimination and uses it only as a rewriting of the same condition (its equation (admission)); that compilation is not formalized, Tarski's theorem not being in Mathlib. The example of root rank $\omega$ in Section~7 is not formalized. Every quantifier of the class is retained: all initial profiles including boundary faces, simultaneous updates, all tie choices, the exact Nash set, and the Euclidean distance. The formal theorems provide no finite recognition algorithm, and the hand proof does not claim one.

\section{Build and audit}

Lean \texttt{v4.34.1}, Mathlib commit \texttt{d13f23b723b8a846827a245b89c10fc7d3f11612}; about 1,700 lines in 8 files. The commands \texttt{lake exe cache get}, \texttt{lake build} and \texttt{lake env lean Verify.lean} were run on a fresh clone on a separate Linux machine. The build reports no errors and no warnings; the sources contain no \texttt{sorry}. Each of the 25 statements audited in \texttt{Verify.lean}, including \texttt{main\_theorem} and \texttt{analyticSet\_compl\_CFP}, depends only on \texttt{propext}, \texttt{Classical.choice} and \texttt{Quot.sound}.

\section{Completion Estimate}
\noindent\textbf{COMPLETION: 100\%}

This is the claim for the full catalogue question as stated by the hand proof: an exact (infinitary) characterization of the games in which fictitious play converges universally to the Nash set, together with the coanalytic bound. The machine check covers all three equivalences of Theorem~\ref{main} and Proposition~\ref{coan}, for every finite game, with no hypothesis and no external theorem. It is not an independent expert review, and the hand proof itself remains an LLM claim pending community review.

\section*{Attribution and reproduction}
Model and process: Claude Fable 5.1 (Claude Code agent session, Anthropic), 11 October 2026, working from the manuscript \cite{Hand} and the Mathlib library; no other model contributed to the Lean development. Human contributions: none; the submitter specified the task and submitted the result, with no mathematical edits. The complete generated output is the Lean repository at its pinned commit \cite{Lean}; the interactive session prompts are not published. Lean source: \url{https://github.com/jbaelaw/fictitious-play-characterization-lean-20261011/tree/d6a956ab7721625c417eb2ffab6e73d5d7a3af49} (\texttt{FPlay/Main.lean}, theorem \texttt{main\_theorem}; \texttt{FPlay/Coanalytic.lean}, theorem \texttt{analyticSet\_compl\_CFP}). Reproduction: Lean \texttt{v4.34.1}, Mathlib commit \texttt{d13f23b723b8a846827a245b89c10fc7d3f11612}, built with \texttt{lake} as described above; the project has several files and depends on Mathlib, so it is not a single-file Lean Web example. Checking outcome: no errors, no warnings, no \texttt{sorry}, standard axioms only, on a fresh clone. The repository maintainers have not verified the result.

\section*{Final status}
SOLVED, as a machine-checked formalization of the hand proof's complete characterization; community expert review of the characterization is pending.

\begin{thebibliography}{9}
\bibitem{Hand} Ji Ho Bae, \emph{An infinitary payoff-array characterization of universal convergence of fictitious play}, hand proof for C73-11, \url{https://github.com/jbaelaw/universal-fictitious-play-characterization/blob/09c4bb2c4b9b7a64c2a9fb87516da03a1b559abc/paper/Proof.tex}.
\bibitem{Lean} Lean 4 formalization, \url{https://github.com/jbaelaw/fictitious-play-characterization-lean-20261011/tree/d6a956ab7721625c417eb2ffab6e73d5d7a3af49}.
\bibitem{Li} Hanyu Li, Wenhan Huang, Zhijian Duan, David Henry Mguni, Kun Shao, Jun Wang and Xiaotie Deng, \emph{A survey on algorithms for Nash equilibria in finite normal-form games}, arXiv:2312.11063 (2023), open problem 5.
\end{thebibliography}
\end{document}
